\documentclass{article}%
\usepackage[T1]{fontenc}%
\usepackage[utf8]{inputenc}%
\usepackage{lmodern}%
\usepackage{textcomp}%
\usepackage{lastpage}%
\usepackage{geometry}%
\geometry{margin=2.2cm}%
%
\usepackage{amsmath}%
\usepackage{amssymb}%
\usepackage{booktabs}%
\usepackage{longtable}%
\usepackage{array}%
\usepackage{hyperref}%
\hypersetup{hidelinks}%
\title{ITA 2027 {-} Química\\[4pt]\large Comparação com o gabarito Poliedro Resolve (ITA 2027, 2ª fase, dia 2 {-} Química)\thanks{Gerado com gpt-resolve: \url{https://github.com/lgabs/gpt-resolve}}}%
\date{29/10/2026}%
%
\begin{document}%
\normalsize%
\maketitle%
\section*{Resumo}\textbf{Nota média: 9,9/10} (36 de 37 itens com a resolução correta; cada item vale o mesmo dentro da questão).\par\medskip%
\begin{longtable}{@{}c c p{0.78\textwidth}@{}}\toprule Questão & Nota & Resumo \\ \midrule
1 & 10,0 & A resolução está certa e bate com o gabarito do Poliedro nos quatro itens. (a) C5H12 + 8 O2 $\to$ 5 CO2 + 6 H2O. (b) 1,6 mol de O2 e 6,4 mol de N2, com n\_ar = 8,0 mol. (c) T\_f = 2300,84 K. (d) Mistura pobre: x $\approx$ 0,1267 mol de pentano, redução de 37\%. Refiz as contas em Python e elas conferem. Nem a resolução nem o gabarito têm erro; só há hipóteses implícitas razoáveis: água como vapor e 1 bar = 10\^{}5 Pa. \\
2 & 10,0 & A resolução da questão 2 está correta e coincide com o gabarito do Poliedro. Os dois dão $\pi$ = 5,008/(10 - x) atm de 0 a 6 dm (de 0,50 a 1,25 atm) e um patamar de $\approx$ 1,25 atm de 6 a 9 dm, por causa da saturação a 0,05 mol/L. Refiz as contas em Python e elas conferem. A resolução declara as hipóteses (i = 1, equilíbrio quase estático, volume do sólido desprezado) e acrescenta a massa precipitada em x = 9 dm (12 g). Não há divergência; ficam só observações de baixa importância sobre hipóteses implícitas do enunciado. \\
3 & 10,0 & Os quatro itens coincidem com o gabarito do Poliedro: X = 2{-}metilbut{-}2{-}eno, Z = 2{-}bromo{-}2{-}metilbutano, G = 2{-}bromo{-}3{-}metilbutano e J = 2{-}metilbutano{-}2,3{-}diol. A química está correta (Saytzeff, Markovnikov, mecanismo radicalar anti{-}Markovnikov e di{-}hidroxilação syn), e a equação de Baeyer está balanceada. As observações de estereoquímica (Z aquiral; G e J racematos) também estão corretas. O único ponto discutível é do enunciado: usar H2O2 aquoso como peróxido, o que não afeta a resposta esperada. \\
4 & 10,0 & Questão 4: a resolução está correta nos quatro itens e bate com o gabarito do Poliedro: a) m = 40 g; b) P = 75\%; c) MM = 180 g/mol; d) {[}X$^{-}${]} = 3 mol/L. A transcrição confere com o enunciado da prova original (pág. 5). Refiz as contas e chego aos mesmos valores. Quantidade total de cátions: 2 mol/L $\times$ 0,5 L = 1 mol. Pela proporção 6:3:1, são 0,6 mol de Q?, 0,3 mol de R$^2$? e 0,1 mol de T$^3$?. Com a pureza de cada sal saem os itens a, b e c. No item d, a quantidade de X$^{-}$ é 1,5 mol, e a eletroneutralidade confirma esse valor. Não encontrei erro em nenhum dos dois. Ficam só duas hipóteses implícitas razoáveis: a leitura da concentração total em mol/L e as impurezas inertes. \\
5 & 10,0 & A resolução está certa e coincide com o gabarito do Poliedro nos três itens. a) A semirreação N2 + 8 H+ + 6 e- $\to$ 2 NH4+ está balanceada. b) Os NOX são +2, +5, 0, -3, +3 e a média +1 no N2O. c) A sequência N2 $\to$ NH4+ $\to$ NO2- $\to$ NO3- $\to$ NO2- $\to$ NO $\to$ N2O $\to$ N2 é a mesma do gabarito. Também conferi o balanceamento das semirreações auxiliares. Sobram duas observações menores: a explicação sobre o N2O cita só uma estrutura de ressonância, e a palavra 'parcial' no enunciado poderia justificar parar em N2O. Nenhuma das duas muda a nota. \\
6 & 10,0 & A resolução está correta. Duas das três reações coincidem com o gabarito: Li + H2O, com coeficientes escalados, e a reforma a vapor, idêntica, com os mesmos estados e condições essencialmente iguais. Para a terceira reação, a resolução escolheu a pirólise do metano e o gabarito, a eletrólise da água. As duas são válidas, porque a questão é aberta. No item iv, a resolução aponta a pirólise (2 mol H2/mol) e o gabarito, a reforma (1,5 mol H2/mol). Refiz as contas e os valores das reações em comum batem (0,25; 1; 1,5). A diferença vem só do conjunto de reações escolhido, e cada um acerta dentro do seu. Numa leitura absoluta, a resposta da resolução é até melhor, pois a pirólise rende mais que a reforma. O principal ponto para um professor é garantir que a correção não exija 'reforma' de forma rígida no item iv. \\
7 & 10,0 & Comparei a resolução da questão 7 com o gabarito do Poliedro e os quatro itens coincidem: os dois estão certos. No item a, as equações das duas etapas (oxidação do etanol a ácido acético e esterificação de Fischer) e as condições de cada uma batem; a equação com dicromato está bem balanceada. Nos itens b, c e d, refiz as contas em Python: $\Delta$H$^\circ$ = -496,2 kJ/mol, $\Delta$S$^\circ$ = -127,2 J/(mol$\cdot$K) e $\Delta$G$^\circ$ $\approx$ -458,3 kJ/mol, que é negativo, então a reação é espontânea a 298 K. Reli o enunciado no original e os dados da transcrição estão corretos. Ficaram só observações menores, nenhuma que mude a nota: a temperatura{-}limite de cerca de 3,9$\times$10$^3$ K, que a resolução acrescenta sem ser pedida, não tem sentido físico, e a resolução fala em \textbackslash{}"meio anidro\textbackslash{}" na esterificação, uma idealização. \\
8 & 10,0 & A resolução está certa nos três itens e confere com o gabarito do Poliedro. O enunciado transcrito corresponde ao original. a) S e Se têm 2 elétrons desemparelhados cada; K tem 1 elétron de valência. b) O S do H2S tem 2 pares não ligantes; o Se do SeF4 tem 1. c) A geometria do H2S é angular; a do SeF4 é gangorra. As configurações eletrônicas e a contagem VSEPR estão corretas. Só há ressalvas pequenas nos comentários extras (a causa do ângulo de \textasciitilde{}92$^\circ$ no H2S e a explicação da expansão do octeto), que não mudam a nota. \\
9 & 10,0 & A resolução está certa e bate com o gabarito do Poliedro nos dois itens: a) {[}HX{]} = 8,0$\cdot$10$^{-}$$^3$ mol/L e b) {[}HY{]} = 1,2$\cdot$10$^{-}$$^2$ mol/L. Os dois seguem o mesmo raciocínio. A soma dos mols dos ácidos vem dos 20 mL de base. A partir do pH 9 no ponto de equivalência (hidrólise de Y$^{-}$ com Kb = 10$^{-}$$^{8}$), sai {[}Y$^{-}${]} = 10$^{-}$$^2$ mol/L em 120 mL. Resolvi o balanço de cargas exato e obtive 7,99$\cdot$10$^{-}$$^3$ e 1,201$\cdot$10$^{-}$$^2$ mol/L, o que confirma as aproximações. O gabarito só tem erros de digitação sem efeito no resultado ({[}HX{]} escrito no lugar de {[}HY{]} na expressão de K\_h). \\
10 & 8,8 & Conferi o enunciado no original (p. 12 do PDF), e a transcrição está correta. Item a: a resolução e o gabarito coincidem nos quatro casos. Em (1), (2) e (3), X é o cátodo e Y é o ânodo; em (4), a relação se inverte, com X (alta {[}Cl-{]}) como ânodo e Y como cátodo. Item b: as duas fontes indicam proteção catódica em (1), (3) e (4). Em (2) elas divergem: a resolução indica anódica e o gabarito, catódica. O gabarito é o mais defensável, porque em pH baixo o ferro não tem domínio passivo e a polarização anódica aceleraria a corrosão da região Y; a própria resolução usa esse argumento em a(2) e depois o contradiz. Resultado: 7 de 8 subitens corretos, com erro em b(2). \\
\bottomrule\end{longtable}%
\section*{Principais pontos de dúvida}%
\begin{itemize}%
\item%
\textbf{Q6} \textit{(key)} No item iv, a resposta do gabarito (reforma a vapor, 1,5 mol H2/mol) não é única. Ela depende das três reações que o candidato escolhe, e a questão é aberta (o próprio gabarito diz 'por exemplo'). Quem escrever a pirólise do CH4 (CH4 $\to$ C + 2 H2, permitida por 'reagindo individualmente') obtém corretamente 2 mol H2/mol, valor acima da reforma. Um corretor que siga o gabarito à risca pode descontar indevidamente o item iv da resolução. Vale confirmar com um professor que a correção deve aceitar a reação de maior rendimento entre as escritas pelo candidato.%
\item%
\textbf{Q10} \textit{(solution)} b(2): a resolução indica proteção anódica para a barra com pH diferencial e o gabarito indica catódica. A escolha da resolução contradiz o próprio argumento de Pourbaix que ela usa no item a(2): em pH baixo, o Fe não tem domínio passivo. Além disso, o caso do 'aço em H2SO4' só vale para ácido concentrado, que o enunciado não especifica. Pelo gabarito, a resolução perde b(2). Vale um professor confirmar se alguma banca aceitaria a resposta 'anódica' com essa justificativa (alto custo de corrente catódica em meio ácido).%
\end{itemize}%
\section*{Questão 1}%
\href{https://poliedroresolve.sistemapoliedro.com.br/vestibulares/ita/2027/ita-2-fase-dia-2-21-10-2026-2-fase-dia-2/dissertativa-1-quimica-2-fase-dia-2-ita-2027}{Resolução do gabarito}\par\smallskip%
\begin{longtable}{@{}p{0.06\textwidth} p{0.33\textwidth} p{0.33\textwidth} p{0.1\textwidth} p{0.08\textwidth}@{}}\toprule Item & Resolução & Gabarito & Veredito & Certo \\ \midrule
a & C5H12 + 8 O2 $\to$ 5 CO2 + 6 H2O ($\Delta$Hc = -3440 kJ/mol) & C5H12 + 8 O2 $\to$ 5 CO2 + 6 H2O (com ar: C5H12 + 8 O2 + 32 N2 $\to$ 5 CO2 + 6 H2O + 32 N2) & igual & ambos \\
b & n\_ar = 8,0 mol; O2 = 1,6 mol; N2 = 6,4 mol & n\_ar = 8 mol; O2 = 1,6 mol; N2 = 6,4 mol & igual & ambos \\
c & 0,20 mol de pentano, 688 kJ, 8,6 mol de gases, $\Delta$T = 2000 K, T\_f = 2300,84 K & 0,2 mol de pentano; $\Delta$T = 2000 K $\Rightarrow$ T\_f = 2300,84 K & igual & ambos \\
d & mistura pobre; x = 416/3284 $\approx$ 0,127 mol; redução de $\approx$ 37\% (fica em \textasciitilde{}63\% do estequiométrico) & reduzir em $\approx$ 37\% (0,1267 mol de pentano, mistura pobre) & igual & ambos \\
\bottomrule\end{longtable}%
\begin{itemize}%
\item%
\textbf{a)} Mesma equação balanceada (C: 5; H: 12 = 6$\cdot$2; O: 16 = 10 + 6). O gabarito traz também a versão com o N2 do ar, mas isso é só um complemento e não muda a resposta.%
\item%
\textbf{b)} Pela conta, n = 1e5$\cdot$0,200/(8,31$\cdot$300,84) = 8,00006 mol. Com 20\% e 80\% em mol, saem 1,6 e 6,4 mol.%
\item%
\textbf{c)} Os produtos somam 1,0 + 1,2 + 6,4 = 8,6 mol. Daí 0,2$\cdot$3440e3/(8,6$\cdot$40) = 2000 K exatos, e T\_f = 2300,84 K.%
\item%
\textbf{d)} O balanço é o mesmo nos dois: 3440e3$\cdot$x = 40$\cdot$(8+3x)$\cdot$1300, que dá x = 0,12667 mol. A redução é (0,2 - 0,12667)/0,2 = 36,67\%, ou seja, 37\% em número inteiro. A resolução ainda confere que sobra O2 ($\approx$ 0,59 mol) e justifica por que descarta a mistura rica.%
\end{itemize}%
\paragraph{Pontos de dúvida}%
\begin{itemize}%
\item%
\textit{(solution)} A resolução e o gabarito tratam a H2O dos produtos como vapor (Cp da mistura gasosa igual a 40). O enunciado só diz isso implicitamente, ao listar H2O entre os gases da mistura. É uma hipótese razoável.%
\item%
\textit{(solution)} Para obter 8,0 mol, a resolução usa 1 bar = 1,0$\times$10\^{}5 Pa e R = 8,31 J/(K$\cdot$mol). É a conversão exata (1 bar = 10\^{}5 Pa) e é coerente com a tabela de constantes. O valor 300,84 K foi escolhido justamente para dar n = 8,0 mol.%
\item%
\textit{(solution)} A resolução escreve o pentano como C5H12(g). A 300,84 K ele seria líquido (ponto de ebulição $\approx$ 309 K), mas no motor é injetado e vaporizado. Isso não muda nada, porque o $\Delta$Hc foi dado pelo enunciado.%
\item%
\textit{(statement)} O item d admite duas soluções: mistura pobre ou mistura rica, esta com cerca de 4,8 mol de pentano em excesso não queimado e supondo Cp = 40 também para o pentano. O critério 'economia de combustível' do enunciado escolhe a mistura pobre, e assim fizeram a resolução e o gabarito. Sobre a mistura rica, a resolução só faz uma observação e não a usa como resposta.%
\end{itemize}%
\section*{Questão 2}%
\href{https://poliedroresolve.sistemapoliedro.com.br/vestibulares/ita/2027/ita-2-fase-dia-2-21-10-2026-2-fase-dia-2/dissertativa-2-quimica-2-fase-dia-2-ita-2027}{Resolução do gabarito}\par\smallskip%
\begin{longtable}{@{}p{0.06\textwidth} p{0.33\textwidth} p{0.33\textwidth} p{0.1\textwidth} p{0.08\textwidth}@{}}\toprule Item & Resolução & Gabarito & Veredito & Certo \\ \midrule
Q2 (gráfico $\pi$ $\times$ deslocamento) & $\pi$(x) = 0,2$\cdot$RT/(10 - x) = 5,008/(10 - x) atm para 0 $\leq$ x $\leq$ 6 dm: ramo de hipérbole crescente de 0,50 a 1,25 atm. Em x = 6 dm (V = 4 L) a solução satura a 0,05 mol/L. Para 6 $\leq$ x $\leq$ 9 dm, $\pi$ fica constante em $\approx$ 1,25 atm, com precipitação; em x = 9 dm há 0,15 mol (12 g) de sólido. Soluto tomado como não eletrólito (i = 1). & $\pi$(x) = 0,2$\cdot$0,0821$\cdot$305/(10 - x) $\approx$ 5,0081/(10 - x) atm para 0 $\leq$ x $\leq$ 6 dm (de 0,50 a 1,25 atm). A partir de x = 6 dm a solução satura (0,05 mol/L) e $\pi$ fica constante em $\approx$ 1,252 atm até 9 dm. O gráfico é uma hipérbole crescente seguida de um patamar. Soluto tomado como não eletrólito. & igual & ambos \\
\bottomrule\end{longtable}%
\begin{itemize}%
\item%
\textbf{Q2 (gráfico $\pi$ $\times$ deslocamento))} Refiz as contas. RT = 0,0821$\cdot$305 = 25,04 atm$\cdot$L/mol e nRT = 5,0081 atm$\cdot$L. A concentração de saturação é 4/80 = 0,05 mol/L, o que dá V = 0,2/0,05 = 4 L, ou seja, x = 6 dm. Os pontos da tabela são 0,50; 0,56; 0,63; 0,72; 0,83; 1,00 e 1,25 atm, iguais nos dois. No patamar, $\pi$ = 0,05$\cdot$25,04 = 1,252 atm. A função por partes, o ponto de quebra, os valores e a forma do gráfico coincidem. A resolução ainda informa a massa precipitada em x = 9 dm (12 g), dado extra e correto. O gabarito, por sua vez, confere explicitamente que a concentração inicial (1,6 g/L) é menor que 4 g/L, ou seja, que a solução começa insaturada. A resolução faz a conferência equivalente pelo volume de saturação (4 L < 10 L).%
\end{itemize}%
\paragraph{Pontos de dúvida}%
\begin{itemize}%
\item%
\textit{(statement)} O enunciado não diz se o soluto é eletrólito. Gabarito e resolução declaram i = 1. Se o soluto se dissociasse, todos os valores de $\pi$ seriam multiplicados por i, mas a forma do gráfico (hipérbole até 6 dm, patamar de 6 a 9 dm) seria a mesma.%
\item%
\textit{(statement)} Ambos desprezam o volume do sólido precipitado (12 g em x = 9 dm), a pressão hidrostática e a não idealidade. Também tomam a solubilidade de 4 g/L como por litro de solução. São aproximações razoáveis para uma solução diluída (0,02--0,05 mol/L).%
\item%
\textit{(statement)} O gráfico começa em $\pi$ $\approx$ 0,50 atm em x = 0, e não em zero. Isso pressupõe que, na posição inicial, já atua a sobrepressão de equilíbrio (processo reversível, quase estático). Gabarito e resolução fazem essa mesma leitura, a mais natural, mas um aluno que começasse o gráfico em 0 atm estaria interpretando o 'estado inicial' de outro modo.%
\end{itemize}%
\section*{Questão 3}%
\href{https://poliedroresolve.sistemapoliedro.com.br/vestibulares/ita/2027/ita-2-fase-dia-2-21-10-2026-2-fase-dia-2/dissertativa-3-quimica-2-fase-dia-2-ita-2027}{Resolução do gabarito}\par\smallskip%
\begin{longtable}{@{}p{0.06\textwidth} p{0.33\textwidth} p{0.33\textwidth} p{0.1\textwidth} p{0.08\textwidth}@{}}\toprule Item & Resolução & Gabarito & Veredito & Certo \\ \midrule
a & X = 2{-}metilbut{-}2{-}eno, (CH3)2C=CHCH3 (produto de Saytzeff da desidratação E1) & X = 2{-}metilbut{-}2{-}eno, (CH3)2C=CH-CH3 & igual & ambos \\
b & Z = 2{-}bromo{-}2{-}metilbutano, (CH3)2CBrCH2CH3 (Markovnikov) & Z = 2{-}bromo{-}2{-}metilbutano, (CH3)2CBr-CH2CH3 (Markovnikov) & igual & ambos \\
c & G = 2{-}bromo{-}3{-}metilbutano, (CH3)2CH-CHBr-CH3 (anti{-}Markovnikov, via radicais) & G = 2{-}bromo{-}3{-}metilbutano, (CH3)2CH-CHBr-CH3 (anti{-}Markovnikov, efeito peróxido) & igual & ambos \\
d & J = 2{-}metilbutano{-}2,3{-}diol, (CH3)2C(OH)-CH(OH)-CH3 (di{-}hidroxilação syn) & J = 2{-}metilbutano{-}2,3{-}diol, (CH3)2C(OH)-CH(OH)-CH3 & igual & ambos \\
\bottomrule\end{longtable}%
\begin{itemize}%
\item%
\textbf{a)} Desidratação E1 do álcool terciário. O carbocátion terciário não sofre rearranjo, e pela regra de Saytzeff o produto principal é o alceno trissubstituído. A resolução cita corretamente o 2{-}metilbut{-}1{-}eno como produto minoritário.%
\item%
\textbf{b)} Adição eletrofílica com carbocátion terciário em C2. A observação de que Z é aquiral está certa, porque o C2 tem dois CH3.%
\item%
\textbf{c)} O Br? se liga ao C3 e gera o radical terciário em C2, que depois abstrai H do HBr. Estrutura e nome IUPAC estão corretos: os localizantes \{2,3\} empatam e a ordem alfabética dá o menor número ao bromo. Também está certo que G é quiral e se forma como racemato.%
\item%
\textbf{d)} Oxidação branda de Baeyer, com éster manganato cíclico e diol vicinal. A equação 3 C5H10 + 2 MnO4- + 4 H2O $\to$ 3 C5H12O2 + 2 MnO2 + 2 OH- está balanceada: C 15/15, H 38/38, O 12/12 e carga -2/-2. Com KMnO4 concentrado e a quente, a clivagem daria propanona e ácido acético, como a resolução diz. J tem um só estereocentro (C3) e sai como racemato.%
\end{itemize}%
\paragraph{Pontos de dúvida}%
\begin{itemize}%
\item%
\textit{(statement)} O enunciado usa H2O2 como 'peróxido' na adição de HBr. Em meio aquoso, o H2O2 não é o iniciador adequado para o efeito Kharasch, e o próprio Poliedro observa que o certo seria um peróxido orgânico ou anidro. Ainda assim, a intenção da banca de cobrar o produto anti{-}Markovnikov é clara, e resolução e gabarito seguem essa leitura.%
\item%
\textit{(statement)} A desidratação produz uma mistura: 2{-}metilbut{-}2{-}eno (principal) e 2{-}metilbut{-}1{-}eno (minoritário). O enunciado pede um único X, o que pressupõe o produto de Saytzeff. A resolução declara essa hipótese e, na observação final, dá os produtos alternativos de G e J caso X fosse o alceno minoritário.%
\end{itemize}%
\section*{Questão 4}%
\href{https://poliedroresolve.sistemapoliedro.com.br/vestibulares/ita/2027/ita-2-fase-dia-2-21-10-2026-2-fase-dia-2/dissertativa-4-quimica-2-fase-dia-2-ita-2027}{Resolução do gabarito}\par\smallskip%
\begin{longtable}{@{}p{0.06\textwidth} p{0.33\textwidth} p{0.33\textwidth} p{0.1\textwidth} p{0.08\textwidth}@{}}\toprule Item & Resolução & Gabarito & Veredito & Certo \\ \midrule
a & m = 40 g & m = 40 g & igual & ambos \\
b & P = 75\% & P = 75\% & igual & ambos \\
c & MM = 180 g/mol & MM = 180 g/mol & igual & ambos \\
d & {[}X$^{-}${]} = 3,0 mol/L & {[}X$^{-}${]} = 3 mol/L & igual & ambos \\
\bottomrule\end{longtable}%
\begin{itemize}%
\item%
\textbf{a)} n\_cátions = 2 mol/L $\cdot$ 0,5 L = 1 mol; n\_Q = 0,6 mol; massa pura = 0,6$\cdot$60 = 36 g; m = 36/0,90 = 40 g. Resolução e gabarito coincidem.%
\item%
\textbf{b)} n\_R = 0,3 mol; massa pura de RX2 = 0,3$\cdot$150 = 45 g; P = 45/60 = 0,75, ou seja, 75\%. Coincidem.%
\item%
\textbf{c)} n\_T = 0,1 mol; massa pura = 0,9$\cdot$20 = 18 g; MM = 18/0,1 = 180 g/mol. Coincidem.%
\item%
\textbf{d)} n\_X$^{-}$ = 0,6 + 2$\cdot$0,3 + 3$\cdot$0,1 = 1,5 mol; {[}X$^{-}${]} = 1,5/0,5 = 3 mol/L. A resolução ainda confere pela eletroneutralidade. Escrever 3,0 em vez de 3 é só notação.%
\end{itemize}%
\paragraph{Pontos de dúvida}%
\begin{itemize}%
\item%
\textit{(statement)} Resolução e gabarito leem "concentração total de cátions = 2 mol/L" como a soma das concentrações em mol/L dos três cátions, e não como concentração em equivalentes de carga. É a leitura natural do enunciado, mas é uma hipótese que nenhum dos dois declara.%
\item%
\textit{(solution)} A resolução supõe que os sais se dissociam totalmente e que as impurezas são inertes, sem Q, R, T nem X. É isso que permite calcular {[}X$^{-}${]} só a partir dos sais puros. A hipótese é razoável e o gabarito também a usa sem dizer.%
\end{itemize}%
\section*{Questão 5}%
\href{https://poliedroresolve.sistemapoliedro.com.br/vestibulares/ita/2027/ita-2-fase-dia-2-21-10-2026-2-fase-dia-2/dissertativa-5-quimica-2-fase-dia-2-ita-2027}{Resolução do gabarito}\par\smallskip%
\begin{longtable}{@{}p{0.06\textwidth} p{0.33\textwidth} p{0.33\textwidth} p{0.1\textwidth} p{0.08\textwidth}@{}}\toprule Item & Resolução & Gabarito & Veredito & Certo \\ \midrule
a & N2 + 8 H+ + 6 e{-} {-}> 2 NH4+ & N2 + 8 H+ + 6 e- $\to$ 2 NH4+ & igual & ambos \\
b & NO: +2; NO3{-}: +5; N2: 0; NH4+: {-}3; NO2{-}: +3; N2O: +1 (médio) & NO: +2; NO3-: +5; N2: 0; NH4+: -3; NO2-: +3; N2O: +1 (médio) & igual & ambos \\
c & N2 $\to$ NH4+ (fixação) $\to$ NO2- $\to$ NO3- (nitrificação) $\to$ NO2- $\to$ NO $\to$ N2O $\to$ N2 (desnitrificação) & N2 $\to$ NH4+ (fixação) $\to$ NO2- $\to$ NO3- (nitrificação) $\to$ NO2- $\to$ NO $\to$ N2O $\to$ N2 (desnitrificação) & igual & ambos \\
\bottomrule\end{longtable}%
\begin{itemize}%
\item%
\textbf{a)} As duas respostas são a mesma equação. Conferi o balanço: N 2=2, H 8=8, carga +8-6=+2=2$\cdot$(+1). O NH4+ é a única espécie da lista com NOX negativo, então é o único produto possível de uma redução do N2.%
\item%
\textbf{b)} Os seis valores são iguais. Refiz cada um somando os NOX até a carga da espécie (H = +1, O = -2) e todos conferem.%
\item%
\textbf{c)} A sequência é idêntica nas duas, com o NO2- aparecendo duas vezes (na nitrificação e na desnitrificação). O NOX fica 0 $\to$ -3 $\to$ +3 $\to$ +5 $\to$ +3 $\to$ +2 $\to$ +1 $\to$ 0, o que bate com o sentido de cada processo. Chequei com script as sete semirreações auxiliares da resolução, e todas estão balanceadas em átomos e carga.%
\end{itemize}%
\paragraph{Pontos de dúvida}%
\begin{itemize}%
\item%
\textit{(solution)} Comentário extra sobre o N2O. A resolução cita só a estrutura N(-)=N(+)=O, em que o N terminal tem NOX -1 e o central +3. A estrutura de ressonância mais importante é N$\equiv$N(+)--O(-), e nela os NOX são 0 e +2. A afirmação não está errada para a estrutura citada e não altera a média +1 pedida no item, mas fica incompleta.%
\item%
\textit{(statement)} O enunciado fala em 'desnitrificação parcial', o que poderia justificar terminar a sequência em N2O em vez de N2. O gabarito e a resolução terminam em N2. A resolução ainda comenta que o processo parcial pode parar em N2O. A ambiguidade não deve afetar a nota.%
\end{itemize}%
\section*{Questão 6}%
\href{https://poliedroresolve.sistemapoliedro.com.br/vestibulares/ita/2027/ita-2-fase-dia-2-21-10-2026-2-fase-dia-2/dissertativa-6-quimica-2-fase-dia-2-ita-2027}{Resolução do gabarito}\par\smallskip%
\begin{longtable}{@{}p{0.06\textwidth} p{0.33\textwidth} p{0.33\textwidth} p{0.1\textwidth} p{0.08\textwidth}@{}}\toprule Item & Resolução & Gabarito & Veredito & Certo \\ \midrule
Reação com Li e H2O (i, ii, iii) & 2 Li(s) + 2 H2O(?) $\to$ 2 LiOH(aq) + H2(g); espontânea em condições ambiente, basta o contato do Li com a água, sem catalisador nem aquecimento & Li(s) + H2O(l) $\to$ LiOH(aq) + 1/2 H2(g); basta o contato direto do lítio com a água & equivalente & ambos \\
Reforma a vapor do metano (i, ii, iii) & CH4(g) + H2O(g) ? CO(g) + 3 H2(g); cerca de 700--1000 $^\circ$C, catalisador de Ni (Ni/Al2O3), excesso de vapor; endotérmica ($\Delta$H$^\circ$ $\approx$ +206 kJ/mol), $\Delta$n\_gás = +2, favorecida por alta T e pressão não muito alta & CH4(g) + H2O(g) $\to$ CO(g) + 3 H2(g); alta temperatura, alta pressão e catalisador adequado & igual & ambos \\
Terceira reação (i, ii, iii) & Pirólise do metano: CH4(g) $\to$ C(s) + 2 H2(g); aquecimento forte (> \textasciitilde{}1000 $^\circ$C sem catalisador, ou 600--900 $^\circ$C com Ni/Fe), sem O2; endotérmica ($\Delta$H$^\circ$ $\approx$ +75 kJ/mol). A eletrólise da água é citada nas verificações como alternativa válida (1 mol H2/mol). & Eletrólise da água: 2 H2O(l) $\to$ 2 H2(g) + O2(g); eletrólito (por exemplo LiOH), eletrodos inertes (grafite ou Pt), corrente contínua & diferente & ambos \\
iv) Reação de maior rendimento de H2 por mol de reagente & Pirólise do metano: 2 mol H2/mol de reagente (contra 1,5 da reforma e 0,25 de Li/H2O; eletrólise = 1). Ressalva: por mol de CH4, a reforma completa CH4 + 2 H2O $\to$ CO2 + 4 H2 daria 4. & Reforma a vapor do metano: 1,5 mol H2/mol de reagente (contra 1 da eletrólise e 0,25 de Li/H2O) & diferente & ambos \\
\bottomrule\end{longtable}%
\begin{itemize}%
\item%
\textbf{Reação com Li e H2O (i, ii, iii))} É a mesma reação, com os mesmos estados físicos. Os coeficientes da resolução são os do gabarito multiplicados por 2. As condições coincidem (contato direto, temperatura ambiente).%
\item%
\textbf{Reforma a vapor do metano (i, ii, iii))} Equação, estados e estequiometria são idênticos. A única diferença está na pressão. O gabarito fala em 'alta pressão', que é a prática industrial (cerca de 20--30 bar, por razões de processo), mas a termodinâmica pede o contrário: $\Delta$n\_gás = +2, então pelo princípio de Le Chatelier a pressão baixa favorece os produtos. A frase da resolução ('pressão não muito alta favorece os produtos') está termodinamicamente correta. O $\Delta$H$^\circ$ $\approx$ +206 kJ/mol citado também está certo.%
\item%
\textbf{Terceira reação (i, ii, iii))} A questão é aberta: pede três reações quaisquer que usem só Li, H2O e CH4, 'individualmente ou combinadas duas a duas', e o gabarito diz 'por exemplo'. A pirólise do CH4 (CH4 reagindo sozinho) atende ao enunciado tanto quanto a eletrólise da H2O (H2O reagindo sozinha). A equação da resolução está balanceada, os estados estão corretos (C sólido) e as condições e o $\Delta$H$^\circ$ $\approx$ +74,8 kJ/mol estão certos. Além disso, a resolução menciona a eletrólise com o mesmo rendimento do gabarito (2/2 = 1).%
\item%
\textbf{iv) Reação de maior rendimento de H2 por mol de reagente)} Refiz as contas usando o mesmo critério dos dois (mol de H2 dividido pela soma dos coeficientes dos reagentes): Li/H2O = 1/4 = 0,25; eletrólise = 2/2 = 1; reforma = 3/2 = 1,5; pirólise = 2/1 = 2; reforma global = 4/3 $\approx$ 1,33. Os valores das reações em comum batem exatamente. A divergência vem só da escolha das reações. Entre as três do gabarito, a reforma (1,5) é a maior; entre as três da resolução, a pirólise (2) é a maior. Os dois estão certos para o conjunto que escolheram. Se o item for lido de forma absoluta ('qual reação deve render mais'), a resolução está até mais completa, porque a pirólise supera a reforma, inclusive na conta por massa (cerca de 25\% contra 18\%). A resolução declara o critério usado e responde de forma clara (pirólise). A ressalva sobre o mol de CH4 é condicional e não cria ambiguidade na resposta principal.%
\end{itemize}%
\paragraph{Pontos de dúvida}%
\begin{itemize}%
\item%
\textit{(key)} No item iv, a resposta do gabarito (reforma a vapor, 1,5 mol H2/mol) não é única. Ela depende das três reações que o candidato escolhe, e a questão é aberta (o próprio gabarito diz 'por exemplo'). Quem escrever a pirólise do CH4 (CH4 $\to$ C + 2 H2, permitida por 'reagindo individualmente') obtém corretamente 2 mol H2/mol, valor acima da reforma. Um corretor que siga o gabarito à risca pode descontar indevidamente o item iv da resolução. Vale confirmar com um professor que a correção deve aceitar a reação de maior rendimento entre as escritas pelo candidato.%
\item%
\textit{(statement)} A expressão 'por mol de reagente' é ambígua: pode significar por mol total de reagentes (soma dos coeficientes) ou por mol de um reagente específico, como o CH4. Gabarito e resolução usam a soma dos coeficientes. A resolução acrescenta que, por mol de CH4, a reforma completa daria 4 mol de H2. A ressalva está correta e não muda a resposta principal.%
\item%
\textit{(key)} O gabarito pede 'alta pressão' para a reforma a vapor. Essa é a prática industrial, mas termodinamicamente a reação ($\Delta$n\_gás = +2) é favorecida por pressão baixa. A formulação da resolução ('pressão não muito alta favorece os produtos') é a termodinamicamente correta. Não deve afetar a nota.%
\item%
\textit{(solution)} Na verificação por massa, a resolução usa M(Li) = 6,94 g/mol, que não aparece na tabela de massas molares da prova. O valor é o real e é usado só numa verificação acessória, que não altera a conclusão (R1 $\approx$ 4\% em massa).%
\item%
\textit{(key)} Na eletrólise, o gabarito sugere LiOH como eletrólito. O LiOH não é uma das substâncias permitidas, mas entra só como eletrólito e não como reagente, então isso não viola o enunciado.%
\end{itemize}%
\section*{Questão 7}%
\href{https://poliedroresolve.sistemapoliedro.com.br/vestibulares/ita/2027/ita-2-fase-dia-2-21-10-2026-2-fase-dia-2/dissertativa-7-quimica-2-fase-dia-2-ita-2027}{Resolução do gabarito}\par\smallskip%
\begin{longtable}{@{}p{0.06\textwidth} p{0.33\textwidth} p{0.33\textwidth} p{0.1\textwidth} p{0.08\textwidth}@{}}\toprule Item & Resolução & Gabarito & Veredito & Certo \\ \midrule
a & Etapa 1: CH3CH2OH + O2 $\to$ CH3COOH + H2O (forma iônica com dicromato: 3 CH3CH2OH + 2 Cr2O7\^{}2{-} + 16 H+ $\to$ 3 CH3COOH + 4 Cr\^{}3+ + 11 H2O), com oxidante forte em excesso (K2Cr2O7 ou KMnO4 em H2SO4) e aquecimento sob refluxo. Etapa 2: esterificação de Fischer CH3COOH + C2H5OH ? CH3COOC2H5 + H2O, com H2SO4 concentrado como catalisador, aquecimento sob refluxo e excesso de etanol e/ou retirada da água. & 1ª etapa: C2H5OH + O2 $\to$ CH3COOH + H2O, com K2Cr2O7/H2SO4 a quente e refluxo; 2ª etapa: CH3COOH + C2H5OH ? CH3COOC2H5 + H2O, com H2SO4 como catalisador e aquecimento sob refluxo. & igual & ambos \\
b & $\Delta$H$^\circ$ = -496,2 kJ/mol & $\Delta$H$^\circ$ = -496,2 kJ/mol & igual & ambos \\
c & $\Delta$S$^\circ$ = -127,2 J/(mol$\cdot$K) & $\Delta$S$^\circ$ = -127,2 J/(K$\cdot$mol) & igual & ambos \\
d & $\Delta$G$^\circ$ $\approx$ -458,3 kJ/mol (reação espontânea a 298 K) & $\Delta$G$^\circ$ $\approx$ -458,3 kJ/mol & igual & ambos \\
\bottomrule\end{longtable}%
\begin{itemize}%
\item%
\textbf{a)} As equações e as condições coincidem com as do gabarito. A resolução vai além dele: dá a equação iônica com dicromato e explica por que o refluxo impede que a oxidação pare no acetaldeído e como deslocar o equilíbrio. Conferi o balanceamento da equação com dicromato (C 6/6, H 34/34, O 17/17, carga +12/+12) e o balanço de elétrons (3$\times$4 e$^{-}$ cedidos = 2$\times$6 e$^{-}$ recebidos): está correto. A soma das duas etapas dá a reação global do enunciado.%
\item%
\textbf{b)} $\Delta$H$^\circ$ = {[}-480,0 + 2(-285,8){]} - {[}2(-277,7) + 0{]} = -1051,6 + 555,4 = -496,2 kJ/mol. Conferido em Python.%
\item%
\textbf{c)} $\Delta$S$^\circ$ = {[}259,4 + 2(69,9){]} - {[}2(160,7) + 205,0{]} = 399,2 - 526,4 = -127,2 J/(mol$\cdot$K). O sinal negativo faz sentido, porque a reação consome 1 mol de gás e só forma líquidos.%
\item%
\textbf{d)} $\Delta$G$^\circ$ = -496,2 - 298(-0,1272) = -496,2 + 37,91 = -458,29 $\approx$ -458,3 kJ/mol. O gabarito escreve só "kJ", sem o "/mol", mas é apenas notação. A conclusão sobre a espontaneidade está correta.%
\end{itemize}%
\paragraph{Pontos de dúvida}%
\begin{itemize}%
\item%
\textit{(solution)} No item d, a resolução acrescenta que a reação seria espontânea até T = $\Delta$H$^\circ$/$\Delta$S$^\circ$ $\approx$ 3,9$\times$10$^3$ K, supondo $\Delta$H$^\circ$ e $\Delta$S$^\circ$ constantes. O enunciado não pede isso, e o valor não tem sentido físico, porque nessa temperatura etanol, acetato de etila e água não seriam líquidos e os dados tabelados não valeriam mais. A conta está certa (496,2/0,1272 $\approx$ 3901 K) e o comentário não afeta a resposta pedida, mas vale indicar como extrapolação apenas formal.%
\item%
\textit{(solution)} No item a, a resolução fala em "meio anidro" na esterificação de Fischer. Isso é uma idealização: na prática se usam ácido acético glacial e etanol com o mínimo de água, mas a reação forma água e o H2SO4 concentrado age como desidratante. Isso não muda a correção do item.%
\end{itemize}%
\section*{Questão 8}%
\href{https://poliedroresolve.sistemapoliedro.com.br/vestibulares/ita/2027/ita-2-fase-dia-2-21-10-2026-2-fase-dia-2/dissertativa-8-quimica-2-fase-dia-2-ita-2027}{Resolução do gabarito}\par\smallskip%
\begin{longtable}{@{}p{0.06\textwidth} p{0.33\textwidth} p{0.33\textwidth} p{0.1\textwidth} p{0.08\textwidth}@{}}\toprule Item & Resolução & Gabarito & Veredito & Certo \\ \midrule
a & S (Z=16): 2 elétrons desemparelhados; Se (Z=34): 2 elétrons desemparelhados; K (Z=19): 1 elétron na camada de valência (4s$^{1}$) & S: 2 desemparelhados; Se: 2 desemparelhados; K: 1 elétron na camada de valência & igual & ambos \\
b & H2S: 2 pares não ligantes no S; SeF4: 1 par não ligante no Se & H2S: 2 pares não ligantes no S; SeF4: 1 par não ligante no Se & igual & ambos \\
c & H2S: angular; SeF4: gangorra (tetraedro distorcido/disfenoidal) & H2S: angular; SeF4: gangorra (seesaw) & igual & ambos \\
\bottomrule\end{longtable}%
\begin{itemize}%
\item%
\textbf{a)} Conferido: S = {[}Ne{]}3s$^2$3p$^{4}$ e Se = {[}Ar{]}3d$^{1}$$^{0}$4s$^2$4p$^{4}$. Nos dois, o np$^{4}$ fica ?? ? ? pela regra de Hund, o que dá 2 elétrons desemparelhados. K = {[}Ar{]}4s$^{1}$, com 1 elétron na camada N.%
\item%
\textbf{b)} No H2S, o S tem 6 elétrons de valência e usa 2 em ligações: (6-2)/2 = 2 pares livres, tipo AX2E2. No SeF4, o Se tem 6 elétrons de valência e usa 4: (6-4)/2 = 1 par livre, tipo AX4E. A contagem total (34 - 8 - 24 = 2 elétrons no Se) confirma o resultado.%
\item%
\textbf{c)} AX2E2 tem arranjo eletrônico tetraédrico e geometria angular. AX4E tem arranjo bipiramidal trigonal com o par livre na posição equatorial, o que dá a gangorra. A justificativa da posição equatorial está certa: lá o par sofre 2 repulsões a 90$^\circ$, contra 3 na posição axial. 'Tetraedro distorcido' e 'disfenoidal' são nomes equivalentes.%
\end{itemize}%
\paragraph{Pontos de dúvida}%
\begin{itemize}%
\item%
\textit{(solution)} A resolução diz que o ângulo H--S--H (\textasciitilde{}92$^\circ$) é menor que 109,5$^\circ$ 'por causa da repulsão dos pares livres'. A VSEPR prevê só um ângulo um pouco abaixo de 109,5$^\circ$ (como os \textasciitilde{}104,5$^\circ$ da H2O). O valor de \textasciitilde{}92$^\circ$ no H2S se explica melhor pela pouca hibridização do S (ligações com caráter quase p puro). É um comentário a mais, que não muda a resposta pedida ('angular').%
\item%
\textit{(solution)} A resolução justifica a expansão do octeto do Se dizendo que ele é do 4.º período. Esse é o argumento comum nos livros didáticos e não vale pontos no item. Hoje a explicação aceita é de hipervalência sem participação relevante de orbitais d, mas a resolução nem chega a citar orbitais d.%
\item%
\textit{(key)} Na resolução do gabarito para o item a, duas das imagens têm nomes de arquivo da questão 3 ('Q03\_A2\_r', 'Q03\_A3\_r'). Pode ser só reaproveitamento de imagens ou um erro de montagem do material do Poliedro. Não foi possível ver as imagens, mas as respostas finais do gabarito estão completas e corretas.%
\end{itemize}%
\section*{Questão 9}%
\href{https://poliedroresolve.sistemapoliedro.com.br/vestibulares/ita/2027/ita-2-fase-dia-2-21-10-2026-2-fase-dia-2/dissertativa-9-quimica-2-fase-dia-2-ita-2027}{Resolução do gabarito}\par\smallskip%
\begin{longtable}{@{}p{0.06\textwidth} p{0.33\textwidth} p{0.33\textwidth} p{0.1\textwidth} p{0.08\textwidth}@{}}\toprule Item & Resolução & Gabarito & Veredito & Certo \\ \midrule
a & {[}HX{]} = 8,0$\cdot$10$^{-}$$^3$ mol/L & {[}HX{]} = 8,0$\cdot$10$^{-}$$^3$ mol/L & igual & ambos \\
b & {[}HY{]} = 1,2$\cdot$10$^{-}$$^2$ mol/L & {[}HY{]} = 1,2$\cdot$10$^{-}$$^2$ mol/L & igual & ambos \\
\bottomrule\end{longtable}%
\begin{itemize}%
\item%
\textbf{a)} Os dois partem de n\_HX + n\_HY = 0,1 $\times$ 0,020 = 2,0$\cdot$10$^{-}$$^3$ mol e descontam n\_HY = 1,2$\cdot$10$^{-}$$^3$ mol, o que dá n\_HX = 8,0$\cdot$10$^{-}$$^{4}$ mol em 0,100 L, ou seja, 8,0$\cdot$10$^{-}$$^3$ mol/L. Resolvi o balanço de cargas exato no ponto de equivalência, sem aproximações, e obtive 7,99$\cdot$10$^{-}$$^3$ mol/L.%
\item%
\textbf{b)} Com Kb = Kw/Ka = 10$^{-}$$^{8}$ e {[}OH$^{-}${]} = {[}HY{]} = 10$^{-}$$^{5}$ mol/L, sai {[}Y$^{-}${]} = 10$^{-}$$^{1}$$^{0}$/10$^{-}$$^{8}$ = 10$^{-}$$^2$ mol/L em 120 mL. Então n\_HY = 1,2$\cdot$10$^{-}$$^3$ mol e {[}HY{]}$_{0}$ = 1,2$\cdot$10$^{-}$$^2$ mol/L. A conta exata, C\_Y = ({[}OH$^{-}${]} - {[}H?{]})(Ka + {[}H?{]})/{[}H?{]}, dá C\_Y = 1,0009$\cdot$10$^{-}$$^2$ mol/L e {[}HY{]}$_{0}$ = 1,201$\cdot$10$^{-}$$^2$ mol/L, o que confirma o valor.%
\end{itemize}%
\paragraph{Pontos de dúvida}%
\begin{itemize}%
\item%
\textit{(statement)} O enunciado não dá a temperatura nem o Kw. A resolução e o gabarito supõem 25 $^\circ$C (Kw = 10$^{-}$$^{1}$$^{4}$) e dizem isso explicitamente. Também fica implícito que os volumes se somam (100 + 20 = 120 mL).%
\item%
\textit{(key)} O gabarito traz erros de digitação: escreve "{[}OH$^{-}${]} = {[}HX{]}" e K\_h = {[}HX{]}{[}OH$^{-}${]}/{[}Y$^{-}${]}, quando o produto da hidrólise de Y$^{-}$ é HY. Também escreve a unidade de {[}Y$^{-}${]} como "mol$\cdot$L", sem o expoente -1. Nenhum desses erros muda o resultado.%
\end{itemize}%
\section*{Questão 10}%
\href{https://poliedroresolve.sistemapoliedro.com.br/vestibulares/ita/2027/ita-2-fase-dia-2-21-10-2026-2-fase-dia-2/dissertativa-10-quimica-2-fase-dia-2-ita-2027}{Resolução do gabarito}\par\smallskip%
\begin{longtable}{@{}p{0.06\textwidth} p{0.33\textwidth} p{0.33\textwidth} p{0.1\textwidth} p{0.08\textwidth}@{}}\toprule Item & Resolução & Gabarito & Veredito & Certo \\ \midrule
a(1) & X ({[}Fe2+{]} alta) é o cátodo e Y ({[}Fe2+{]} baixa) é o ânodo. Pela equação de Nernst, E = E$^\circ$ + 0,0295$\cdot$log{[}Fe2+{]}, então E\_X > E\_Y. & X catódica e Y anódica. O potencial de Fe2+/Fe é maior onde a {[}Fe2+{]} é maior. & igual & ambos \\
a(2) & X (pH alto) é o cátodo e Y (pH baixo) é o ânodo. Em meio ácido o ferro se dissolve ativamente; em meio alcalino ele se passiva e ali ocorre a redução de O2 ou de H2O. & X (pH alto, OH-) catódica e Y (pH baixo, H+) anódica. & igual & ambos \\
a(3) & X (mais O2) é o cátodo e Y (menos O2) é o ânodo. Pela equação de Nernst, E\_X - E\_Y = (0,059/4)$\cdot$log({[}O2{]}\_X/{[}O2{]}\_Y) > 0. & X (O2 alto) catódica e Y (O2 baixo) anódica, por aeração diferencial. & igual & ambos \\
a(4) & X (alta {[}Cl-{]}) é o ânodo, com pites, e Y (baixa {[}Cl-{]}) é o cátodo. & X (Cl- alto) anódica e Y catódica, por pite causado por cloreto. & igual & ambos \\
b(1) & Proteção catódica. & Proteção catódica. & igual & ambos \\
b(2) & Proteção anódica. A resolução argumenta que o ferro já se passiva em meio alcalino, que a proteção anódica estenderia a passivação a toda a barra e que a proteção catódica exigiria uma corrente enorme na região ácida, com risco de fragilização por hidrogênio. & Proteção catódica, como nos demais casos, porque a película passiva do ferro é frágil. & diferente & gabarito \\
b(3) & Proteção catódica, por ser o caso típico de estruturas enterradas ou submersas. & Proteção catódica. & igual & ambos \\
b(4) & Proteção catódica. O Cl- impede a passivação, então a proteção anódica causaria pites. & Proteção catódica. & igual & ambos \\
\bottomrule\end{longtable}%
\begin{itemize}%
\item%
\textbf{a(1))} É a mesma conclusão, com a mesma justificativa: na pilha de concentração iônica, o ânodo é a região em contato com a solução mais diluída. Conferi o fator de Nernst: 2,3RT/F = 0,0590 V.%
\item%
\textbf{a(2))} A conclusão é a mesma. A justificativa da resolução (passivação em meio alcalino e dissolução ativa em meio ácido, pelo diagrama de Pourbaix) é até mais clara que a do gabarito, que fala em 'OH- desfavorece a redução de H+'.%
\item%
\textbf{a(3))} É o caso clássico de aeração diferencial, e as duas fontes concluem o mesmo.%
\item%
\textbf{a(4))} A conclusão é a mesma, e as duas fontes atribuem o efeito ao rompimento da película passiva pelo Cl-. A resolução acrescenta a complexação do Fe2+ como fator que abaixa o potencial local. Isso é um argumento extra, que não muda a resposta.%
\item%
\textbf{b(1))} As duas respostas coincidem.%
\item%
\textbf{b(2))} Refazendo o raciocínio: a proteção anódica só funciona se TODA a superfície puder ser levada à região passiva. Na região Y (pH baixo), o diagrama de Pourbaix do Fe não tem domínio de passivação: em qualquer potencial, as espécies estáveis são Fe2+ ou Fe3+, ambas solúveis. Portanto, polarizar Y anodicamente só levaria o ferro a Fe3+ e aceleraria a corrosão justamente na região que já é anódica. A própria resolução usa esse argumento em a(2) ('o ferro fica na região de corrosão do diagrama de Pourbaix') e depois o contradiz em b(2). O exemplo citado de aço em H2SO4 vale só para ácido sulfúrico concentrado, em que a passivação é cinética, e o enunciado não diz qual é o ácido. O alto custo de corrente da proteção catódica em meio ácido é uma objeção real, mas não a torna inviável, enquanto a proteção anódica pode simplesmente falhar. A resposta do gabarito (catódica) é a mais defensável, embora sua justificativa ('película frágil') seja simplista.%
\item%
\textbf{b(3))} As duas respostas coincidem.%
\item%
\textbf{b(4))} As duas respostas coincidem, e a justificativa da resolução, que contraindica a proteção anódica na presença de cloreto, está correta.%
\end{itemize}%
\paragraph{Pontos de dúvida}%
\begin{itemize}%
\item%
\textit{(solution)} b(2): a resolução indica proteção anódica para a barra com pH diferencial e o gabarito indica catódica. A escolha da resolução contradiz o próprio argumento de Pourbaix que ela usa no item a(2): em pH baixo, o Fe não tem domínio passivo. Além disso, o caso do 'aço em H2SO4' só vale para ácido concentrado, que o enunciado não especifica. Pelo gabarito, a resolução perde b(2). Vale um professor confirmar se alguma banca aceitaria a resposta 'anódica' com essa justificativa (alto custo de corrente catódica em meio ácido).%
\item%
\textit{(key)} A justificativa do gabarito para b ('a película passiva do ferro é frágil, então a proteção anódica nunca é adequada') é simplista. O aço{-}carbono recebe proteção anódica na indústria, por exemplo em H2SO4 concentrado e em soluções alcalinas. Uma justificativa melhor seria caso a caso: (4) o Cl- rompe a película; (2) não há domínio passivo em meio ácido; (1) e (3) são meios neutros sem agente passivante, em que a proteção catódica é o padrão.%
\item%
\textit{(statement)} O enunciado ('para cada caso, itens 1 a 4') sugere que as respostas poderiam variar de caso para caso, mas não especifica o ácido nem os contraíons em (1), (2) e (4). Essa omissão abre espaço para respostas diferentes em b(2), dependendo da hipótese sobre o meio.%
\item%
\textit{(key)} Em a(2), o raciocínio do gabarito ('a presença de OH- desfavorece a redução de H+; assim o ferro oxida menos') é confuso, porque a reação catódica ocorre justamente em X, onde há OH-. A conclusão (X catódica, Y anódica) está correta. A explicação da resolução (passivação em pH alto e dissolução ativa em pH baixo) é mais coerente.%
\item%
\textit{(solution)} Em a(1), a resolução põe a deposição de Fe2+ como reação catódica em X. Essa é a leitura de livro{-}texto para a pilha de concentração iônica, e o gabarito faz o mesmo. Em meio aerado, a redução de O2 também poderia ocorrer. A hipótese é razoável e coincide com o gabarito.%
\end{itemize}%
\end{document}